\subsection{Модуль lab2}

\subsubsection{BFS: обход рёбер}

\textbf{Вход:} \texttt{Graph const\&} (конструктор), \texttt{int start}.

\textbf{Выход:} \texttt{BFSResult}~--- \texttt{visited\_order}
(\texttt{vector<int>}), \texttt{edges\_traversed}
(\texttt{vector<pair<int,int>>}), \texttt{parent}
(\texttt{unordered\_map<int,int>}), \\ \texttt{dist\_steps}
(\texttt{unordered\_map<int,int>}), \texttt{iterations} (\texttt{size\_t}).

\textbf{Описание:} BFS с очередью. Счётчик
\texttt{iterations} инкрементируется при каждом рассмотрении ребра $(u, v)$ и его извлечения из очереди. \texttt{dist\_steps[v] = -1}
означает недостижимость.

\begin{lstlisting}[style=cppstyle, caption={Обход рёбер в ширину (BFS.cc)}]
BFSResult BFS::traverse(int start) const {
    BFSResult result;
    for (int v : m_graph_.vertexIds()) {
        result.dist_steps[v] = -1;
        result.parent[v]     = -1;
    }
    result.dist_steps[start] = 0;
    result.visited_order.push_back(start);

    std::queue<int> q;
    q.push(start);
    while (!q.empty()) {
        int u = q.front(); q.pop();
        ++result.iterations;
        for (auto const& [v, w] : m_graph_.neighbors(u)) {
            ++result.iterations;
            if (result.dist_steps[v] == -1) {
                result.dist_steps[v] = result.dist_steps[u] + 1;
                result.parent[v]     = u;
                result.edges_traversed.emplace_back(u, v);
                result.visited_order.push_back(v);
                q.push(v);
            }
        }
    }
    return result;
}
\end{lstlisting}

\subsubsection{Dijkstra: кратчайшие пути}

\textbf{Вход:} \texttt{Graph const\&} (конструктор), \texttt{int start}.

\textbf{Выход:} \texttt{DijkstraResult}~--- \texttt{dist}
(\texttt{unordered\_map<int,double>}), \texttt{parent}
(\texttt{unordered\_map<int,int>}), \texttt{iterations} (\texttt{size\_t}).

\textbf{Описание:} Min-heap на \texttt{priority\_queue} с
\texttt{std::greater<>}. Счётчик \texttt{iterations} инкрементируется
при каждой попытке релаксации ребра и работой с кучей. Устаревшие записи в
куче отфильтровываются проверкой \texttt{d > dist[u]}.
Путь восстанавливается через \texttt{PathUtils::reconstructPath}
по массиву \texttt{parent}. Сложность: $O((V+E)\log{V})$.

\begin{lstlisting}[style=cppstyle, caption={Алгоритм Дейкстры (Dijkstra.cc)}]
DijkstraResult Dijkstra::compute(int start) const {
    DijkstraResult result;
    for (int v : m_graph_.vertexIds()) {
        result.dist[v]   = INF;
        result.parent[v] = -1;
    }
    result.dist[start] = 0.0;

    using Entry = std::pair<double, int>;
    std::priority_queue<Entry, std::vector<Entry>,
                        std::greater<Entry>> pq;
    pq.push({0.0, start});

    while (!pq.empty()) {
        auto [d, u] = pq.top(); pq.pop();
        ++result.iterations;
        if (d > result.dist[u]) continue;
        for (auto const& [v, w] : m_graph_.neighbors(u)) {
            ++result.iterations;
            double nd = result.dist[u] + w;
            if (nd < result.dist[v]) {
                result.dist[v]   = nd;
                result.parent[v] = u;
                pq.push({nd, v});
            }
        }
    }
    return result;
}
\end{lstlisting}

\subsubsection{Comparator: сравнение алгоритмов}

\textbf{Вход:} \texttt{Graph const\&}, \texttt{int start}.

\textbf{Выход:} \texttt{CompareResults}~$=$~\texttt{vector<CompareResult>},
каждый элемент содержит \texttt{algorithm\_name} и \texttt{iterations}.

\textbf{Описание:} Запускает BFS и Дейкстру на одном графе с одной
стартовой вершиной и собирает счётчики итераций в единый вектор.

\begin{lstlisting}[style=cppstyle, caption={Сравнение алгоритмов (Comparator.cc)}]
Comparator::CompareResults
Comparator::compare(Graph const& graph, int start) {
    CompareResults results;
    {
        BFS bfs(graph);
        auto r = bfs.traverse(start);
        results.push_back({"BFS", r.iterations});
    }
    {
        Dijkstra dijk(graph);
        auto r = dijk.compute(start);
        results.push_back({"Dijkstra", r.iterations});
    }
    return results;
}
\end{lstlisting}
